% 有限随机矩阵作为概率核。教学数据：输入 S={1,2,3}，输出 T={a,b,c}，
% K 的三行依次为 (0.7,0.2,0.1)、(0.2,0.6,0.2)、(0.1,0.3,0.6)，每行和为 1。
\begin{tikzpicture}[line width=.6pt,font=\fontsize{8}{10}\selectfont,>=Stealth,
  st/.style={circle,draw=FigureBlue,fill=FigureBlueFill,minimum size=6mm,inner sep=0pt},
  ot/.style={circle,draw=FigureGreen,fill=FigureGreenFill,minimum size=6mm,inner sep=0pt},
  cell/.style={draw=FigureStroke,minimum width=10mm,minimum height=6mm,inner sep=0pt},
  note/.style={FigureMuted,font=\fontsize{7.5}{9}\selectfont}]
% ---------- (a) 转移图 ----------
\begin{scope}
\foreach \s/\y in {1/0,2/-12,3/-24}{\node[st] (s\s) at (0,\y mm) {$\s$};}
\foreach \t/\y in {a/0,b/-12,c/-24}{\node[ot] (t\t) at (34mm,\y mm) {$\t$};}
\node[FigureBlue] at (0,8mm) {输入$s$};
\node[FigureGreen] at (34mm,8mm) {输出$t$};
\foreach \s/\t in {1/a,1/b,1/c,3/a,3/b,3/c}{\draw[->,FigureMuted!45,line width=.4pt] (s\s) -- (t\t);}
\draw[->,FigureBlue,line width=.6pt] (s2) -- node[pos=.62,above left=-1pt,fill=white,inner sep=1pt]{$0.2$} (ta);
\draw[->,FigureBlue,line width=1.3pt] (s2) -- node[pos=.62,above,fill=white,inner sep=1pt]{$0.6$} (tb);
\draw[->,FigureBlue,line width=.6pt] (s2) -- node[pos=.62,below left=-1pt,fill=white,inner sep=1pt]{$0.2$} (tc);
\node at (17mm,-34mm) {(a) 输入$2$发出的概率};
\end{scope}
% ---------- (b) 随机矩阵 ----------
\begin{scope}[xshift=60mm]
\foreach \t/\j in {a/0,b/1,c/2}{\node[FigureGreen] at (10mm*\j+5mm,8mm) {$\t$};}
\node[FigureInk] at (37mm,8mm) {行和};
\foreach \s/\i/\a/\b/\c in {1/0/0.7/0.2/0.1,2/1/0.2/0.6/0.2,3/2/0.1/0.3/0.6}{
  \node[FigureBlue,anchor=east] at (-1.5mm,-12mm*\i) {$K(\s,\cdot)$};
  \node[cell,fill=FigureBlueFill] at (5mm,-12mm*\i) {$\a$};
  \node[cell,fill=FigureBlueFill] at (15mm,-12mm*\i) {$\b$};
  \node[cell,fill=FigureBlueFill] at (25mm,-12mm*\i) {$\c$};
  \node[cell,fill=FigureYellowFill] at (37mm,-12mm*\i) {$1$};}
\node at (18mm,-34mm) {(b) 同一核写成随机矩阵};
\end{scope}
\end{tikzpicture}
